\documentclass[10pt,a4paper]{amsart}
\usepackage[margin=19mm]{geometry}
\usepackage{amsmath,amssymb,mathtools}
\usepackage[scr=rsfs,cal=euler]{mathalfa}
\usepackage{xfrac}
\usepackage[unicode,hidelinks,bookmarks=false]{hyperref}
\newcommand{\ie}{i.e.\ }
\newcommand{\cf}{cf.\ }
\newcommand{\resp}{respectively\ }
\newcommand{\Subsection}{\subsection}
\providecommand{\nobreakdash}{\nobreak\hbox{-}}
\setlength{\emergencystretch}{2em}
\multlinegap0pt
\usepackage{bm}
\usepackage{bbm}
\usepackage{dsfont}
\usepackage{xspace}
\usepackage{cancel}
\usepackage{mathtools}
\usepackage{amsthm}
\makeatletter
\@ifundefined{thm}{\newtheorem{thm}{Theorem}}{}
\@ifundefined{lem}{\newtheorem{lem}[thm]{Lemma}}{}
\makeatother
\title[CR Yamabe Equation on the Heisenberg Group via the method of]{CR Yamabe Equation on the Heisenberg Group via the method of moving spheres}
\author{Congwen Liu}
\date{Source: arXiv:2512.22458v1 (CC BY 4.0)}
\begin{document}
\maketitle

\noindent\emph{Source and licence.} CR Yamabe Equation on the Heisenberg Group via the method of moving spheres. Version arXiv:2512.22458v1 (CC BY 4.0), distributed under the Creative Commons Attribution 4.0 International licence, as recorded in the arXiv licence metadata for this exact version and shown on the abstract page https://arxiv.org/abs/2512.22458v1. Full rights notice: licenses/arxiv-2512.22458-CC-BY-4.0.txt.\par\smallskip

\noindent\textit{Editorial scope.} Counted unit: the Lem whose source label is \texttt{lem:asymp} in the article (source0.tex lines 662--674, source label \texttt{lem:asymp}), delivered together with the article's own complete original proof of it (source0.tex lines 676--780, one \texttt{proof} environment of 105 lines). No mathematics is written, repaired or re-derived: statement and proof are the article's own text line for line. Required context retained uncounted: thm:lnzlem2 (lines 217--232); eqn:alphabeta (lines 225--227); eqn:alphaf1 (lines 602--604); lem:fixedpoint1 (lines 629--636). Every definition, display and label of those lines is retained unchanged. Omitted from this package and not redistributed: 1--216, 228--601, 605--628, 637--661, 675--675, 781--1478. Nothing inside a retained excerpt is omitted or altered: every retained body is delivered exactly as the article's lines are written, so no omission receipt of any kind accompanies this package. Citations: the retained excerpts carry no citation command at all, so no bibliography entry of the article is transcribed and no reference is redistributed. Reference adaptation: none was needed and none was made; every reference inside the delivered unit resolves inside it, so no unresolved reference or citation is produced. Printed slips kept literal: no printed word, symbol, hypothesis or number of the counted statement or of its proof was altered. Layout and class: the article's own class is \texttt{amsart} with packages including amscd, graphics, graphicx, subfigure, float, caption2, color, epstopdf, amssymb, mathrsfs, amsmath, amsthm, bbm, cases; the wrapper is the lane's pinned \texttt{amsart} editorial wrapper at 10pt on A4, which restates the theorem-like environments the retained text needs and adds the article's own macro definitions that the retained text needs (none), with extra wrapper packages (bm, bbm, dsfont, xspace, cancel, mathtools). The delivered numbering is the unit's own. Assets: no retained span contains a figure environment, a graphics-inclusion command, a PDF-page inclusion, a table or a TikZ picture; no image file is copied into this package, the package declares no asset dependency and no image file is redistributed with it. Licence: version arXiv:2512.22458v1 of this article is distributed under the Creative Commons Attribution 4.0 International licence, as recorded in the arXiv licence metadata for this exact version and shown on the abstract page https://arxiv.org/abs/2512.22458v1. A full rights notice is delivered at licenses/arxiv-2512.22458-CC-BY-4.0.txt. Characters: every retained excerpt is pure ASCII, so the delivered engine, XeLaTeX with its default Latin Modern text font, reports no missing character.
\begin{thm}[Calculus Lemma II of Li-Nirenberg-Zhu type] \label{thm:lnzlem2}
Let \(n\ge 1\) and \(\nu>0\). Suppose a nonnegative function \(f\in C^0(\mathbb{H}^n)\) attains its maximum at the origin, and for every \(\xi\in\mathbb{H}^n\) there exists \(\lambda(\xi)>0\) such that
\[
\Bigl(\frac{\lambda(\xi)}{d_H(\xi,\zeta)}\Bigr)^{\nu}
f\!\bigl(\Phi_{\xi,\lambda(\xi)}^{\beta_f}(\zeta)\bigr)=f(\zeta)
\qquad \forall \zeta\in\mathbb{H}^n\setminus\{\xi\},
\]
where 
\begin{equation}\label{eqn:alphabeta}
\beta_f:=(\alpha_f)^{2/\nu}f(0)^{-2/\nu} \quad \text{and} \quad \alpha_f:=\lim_{|\zeta|_H\to\infty}|\zeta|_H^{\nu}f(\zeta).
\end{equation}
Then
\[
f(z,t) = \alpha_f\,\bigl| t + i|z|^2 + i\beta_f \bigr|^{-\frac{\nu}{2}}.
\]
\end{thm}

\begin{equation}\label{eqn:alphaf1}
\lambda(\xi)^{\nu} f(\xi) = \alpha_f, \quad \forall \xi \in \mathbb{H}^{n} . 
\end{equation}

\begin{lem}\label{lem:fixedpoint1}
Let $f$ be the function from Theorem~\ref{thm:lnzlem2} and $\alpha_f$, $\beta_f$ be defined as in \eqref{eqn:alphabeta}.
For any $\epsilon>0$, there exists $C_{\epsilon}$ such that for any $|\zeta|_{\! H} \geq C_{\epsilon}$, there exists a point $\xi^{\ast}=\xi^{\ast}(\zeta) \in \mathbb{H}^n$ satisfying
\begin{equation}\label{fixedpoint1}
\Phi_{\xi^{\ast}, \lambda(\xi^{\ast})}^{\beta_f}(\zeta) = 0 
\quad \text { and } \quad \big|\xi^{\ast}\big|_{\! H} \leq \epsilon. 
\end{equation}
\end{lem}

\begin{lem}\label{lem:asymp}
Let $f$ be the function from Theorem~\ref{thm:lnzlem2}. Suppose for small $\delta>0$ the maps $z \colon (-\delta, \delta) \setminus \{0\} \to \mathbb{C}^n$ and $t \colon (-\delta, \delta) \setminus \{0\} \to \mathbb{R}$ are continuous and satisfy one of the following two asymptotic conditions:
\begin{enumerate}
\item[(A1)]\, $z(h) = O(1)$ and $|t(h)| \approx |h|^{-1}$;
\item[(A2)]\, $t(h) = O(1)$ and $|z(h)| \approx |h|^{-1}$.
\end{enumerate}
Here, $a(h) \approx b(h)$ means that the ratio $a(h)/b(h)$ is bounded above and below by positive constants as $h \to 0$.

Then 
\begin{equation}\label{eqn:asym2}
\lim_{h \to 0} \, \frac{1}{h} \Big\{ \bigl|(z(h), t(h))\bigr|_{\! H}^{\nu}  f(z(h), t(h)) - \alpha_f \Big\} = 0.
\end{equation}
\end{lem}

\begin{proof}
Suppose that Condition (A1) is satisfied. It is then straightforward to verify that for any $(z', t') \in \mathbb{H}^n$ and sufficiently small $|h|$, 
\[
\bigl|(z(h),t(h)) \bigr|_{\! H}^4 \sim |t(h)|^2,
\]
and
\[
d_{\! H} \!\left((z(h), t(h)), (z', t')\right)^4 - \bigl|(z(h),t(h)) \bigr|_{\! H}^4
= \bigl[ -2t' + 4 \operatorname{Im} \langle z(h), z' \rangle \bigr] t(h) + o(t(h)).
\]
Here and throughout, $a(h) \sim b(h)$ means that $\lim_{h\to 0} a(h)/b(h) =1$. 
This should cause no confusion with the notation $a(h) \approx b(h)$.
It follows that
\begin{align}
& \bigl|(z(h), t(h)) \bigr|_{\! H}^{\nu} f(z(h), t(h)) \notag \\
&\quad = \lambda(z', t')^{\nu} \left\{ \frac{\bigl|(z(h),t(h)) \bigr|_{\! H}^4}{d_{\! H} \!\left((z(h), t(h)), (z', t')\right)^4} \right\}^{\nu/4} 
f \!\left( \Phi_{(z', t'), \lambda(z', t')}^{\, \beta_f} (z(h), t(h)) \right) \notag \\
&\quad = \lambda(z', t')^{\nu} \left\{ 1 + \frac{\nu}{2} \cdot \frac{t' - 2 \operatorname{Im} \langle z(h), z' \rangle}{t(h)} + o(h) \right\}
f \!\left( \Phi_{(z', t'), \lambda(z', t')}^{\, \beta_f} (z(h), t(h)) \right). \notag
\end{align}
Together with \eqref{eqn:alphaf1}, this implies
\begin{align}
& \frac{1}{|h|} \Bigl\{ \bigl|(z(h),t(h)) \bigr|_{\! H}^{\nu} f(z(h), t(h)) - \alpha_f \Bigr\} \notag \\
&\quad = \lambda(z', t')^{\nu} \cdot \frac{ f \!\left( \Phi_{(z', t'), \lambda(z', t')}^{\, \beta_f}
(z(h), t(h)) \right) - f(z', t') }{|h|} \notag \\
&\qquad + \lambda(z', t')^{\nu} \left\{ \frac{\nu}{2} \cdot \frac{t' - 2 \operatorname{Im} \langle z(h), z' \rangle}{t(h) |h|} + o(1) \right\}
f \!\left( \Phi_{(z', t'), \lambda(z', t')}^{\, \beta_f} (z(h), t(h)) \right). \label{eqn:asym02}
\end{align}
Taking $(z', t') = (0, 0)$ in the above and noting that $f$ attains its maximum at the origin, we obtain
\begin{equation}
\limsup_{h \to 0} \, \frac{1}{|h|} \Bigl\{ \bigl|(z(h),t(h)) \bigr|_{\! H}^{\nu} f(z(h), t(h)) - \alpha_f 
\Bigr\} ~\leq~ 0.
\end{equation}

On the other hand, Condition (A1) implies that there exist constants $c_1, c_2 > 0$ such that $|z(h)| \leq c_1$ and $|t(h) h| \geq c_2$ for all $h \in (-\delta, \delta) \setminus \{0\}$. Hence, for any $\epsilon \in (0, 1)$,
\[
\max_{|(z', t')|_{\! H} \leq \epsilon} \left| \frac{t' - 2 \operatorname{Im} \langle z(h), z' \rangle}{t(h) |h|} \right|
\leq \frac{1 + 2c_1}{c_2} \, \epsilon.
\]
Now, letting $(z', t')$ be the point $\xi^*$ from Lemma~\ref{lem:fixedpoint1}, it follows from \eqref{eqn:asym02} that
\begin{equation*}
\liminf_{h \to 0} \, \frac{1}{|h|} \Bigl\{ \bigl|(z(h),t(h)) \bigr|_{\! H}^{\nu} f(z(h), t(h)) - \alpha_f \Bigr\}
\geq - \epsilon \cdot \frac{\nu (1 + 2c_1) \alpha_f}{2 c_2} f(0) \max_{|\xi|_{\! H} \leq \epsilon} f(\xi)^{-1}.
\end{equation*}
Letting $\epsilon \to 0$, we conclude that
\begin{equation}
\liminf_{h \to 0} \, \frac{1}{|h|} \Bigl\{ \bigl|(z(h),t(h)) \bigr|_{\! H}^{\nu} f(z(h), t(h)) 
- \alpha_f \Bigr\} ~\geq~ 0.
\end{equation}
Therefore, \eqref{eqn:asym2} is established.

We now assume that Condition (A2) holds. In this case, for any $(z',t') \in \mathbb{H}^{n}$
and small $|h|$, we have
\[
\bigl|(z(h),t(h)) \bigr|_{\! H} \sim |z(h)|
\]
and
\[
d_{\! H}  \bigl((z(h), t(h)), (z', t')\bigr)^4 - \bigl|(z(h), t(h)\bigr|_{\! H}^4
= - 4 |z(h)|^2 \operatorname{Re} \langle z(h), z'\rangle 
+ o(|z(h)|^3).
\]
It follows that
\begin{align}
& \bigl|(z(h),t(h)) \bigr|_{\! H}^{\nu} f(z(h),t(h)) \notag\\
&\quad =~ \lambda(z',t')^{\nu} \left\{ \frac {\bigl|(z(h),t(h)) \bigr|_{\! H}^4} 
{d_{\! H} \! \left((z(h), t(h)), \left(z', t'\right)\right)^4}\right\}^{\nu/4} 
f\left(\Phi_{(z',t'),\lambda(z',t')}^{\, \beta_f} (z(h),t(h))\right) \notag\\
&\quad =~ \lambda(z',t')^{\nu} \left\{ 1+ \nu\, 
\frac {\operatorname{Re} \langle z(h), z'\rangle} {|z(h)|^2}
+ o \left( h \right) \right\} 
f\left(\Phi_{(z',t'),\lambda(z',t')}^{\, \beta_f} (z(h),t(h))\right), \notag
\end{align}
hence
\begin{align}
&\frac {1}{|h|} \, \big\{ \bigl|(z(h),t(h)) \bigr|_{\! H}^{\nu} f(z(h),t(h)) - \alpha_f \big\}\notag\\
& \quad \; ~=~ \lambda(z',t')^{\nu} \, \frac {f\left(\Phi_{(z', t'), \lambda(z',t')}^{\, \beta_f}
(z(h),t(h))\right) -f(z',t')}{|h|} \notag\\
& \qquad \quad + \lambda(z',t')^{\nu} \left\{  \nu\, 
\frac {\operatorname{Re} \langle z(h), z'\rangle} {|z(h)|^2 |h|}  + o(1) \right\} \, f\left(\Phi_{(z',t'),\lambda(z',t')}^{\, \beta_f} (z(h),t(h))\right).  \label{eqn:asym01}
\end{align}
Again, taking $(z',t')=(0,0)$ in the above and using the fact that $f$ has a maximum point at the origin, we obtain
\begin{equation}\label{eqn:limsup2}
\limsup_{h\to 0}\, \frac {1}{|h|} \, \Big\{ \bigl|(z(h),t(h))\bigr|_{\! H}^{\nu} f(z(h),t(h)) - \alpha_f \Big\} ~\leq~ 0.
\end{equation}

On the other hand, Condition (A2) implies that there exists a constant $c'_2>0$ such 
that $|z(h)||h|\geq c_2'$ for all $h\in (-\delta,\delta)\setminus \{0\}$. Thus,
for any $\epsilon\in (0,1)$,  
\[
\max_{|(z',t')|_{\! H} \leq \epsilon} \left| \frac {\operatorname{Re}\, 
\langle z(h), z'\rangle}  {|z(h)|^2 |h|} \right| ~\leq~  \frac {1}{c_2'} \epsilon. 
\]
Letting $(z',t')$ be the point $\xi^{\ast}$ from Lemma \ref{lem:fixedpoint1}, it follows from \eqref{eqn:asym01} that
\begin{equation*}
\liminf_{h\to 0}\, \frac {1}{|h|} \, \Big\{ \bigl|(z(h),t(h)) \bigr|_{\! H}^{\nu} f(z(h),t(h)) - \alpha_f \Big\} 
~\geq~ - \epsilon\, \frac {\nu \alpha_f}{c_2'}  f(0) \max_{|\xi|_{\! H} \leq \epsilon} f(\xi)^{-1}.
\end{equation*}
Sending $\epsilon$ to 0 , we have
\begin{equation}
\liminf_{h\to 0}  \frac {1}{|h|} \, \Big\{ \bigl|(z(h),t(h)) \bigr|_{\! H}^{\nu} f(z(h),t(h)) - \alpha_f \Big\}
~\geq~ 0.
\end{equation}
Together with \eqref{eqn:limsup2}, this completes the proof of the lemma.
\end{proof}

\end{document}
